\section{Introduction}\label{introduction}

\subsection{Motivation and relevance}\label{motivation-and-relevance}

Almaty's winter air pollution is one of the most severe urban public-health problems in Central Asia, and its physical origins are well understood. The city sits in a natural mountain basin at an elevation of 700--900 m above sea level, enclosed to the south by the Zailiyskiy Alatau range and open to the steppe to the north. Under anticyclonic winter conditions, a stable temperature inversion forms over the basin, suppressing vertical mixing and trapping combustion products near the surface. During the most intense episodes the atmospheric boundary-layer height (BLH) collapses from a summer-season mean of around 1,000 m to values of 10--50 m, reducing the effective volume available for dilution by a factor of 20 to 100. The primary emission source during these episodes is private-sector coal combustion: the great majority of residential heating in peri-urban Almaty relies on solid-fuel boilers, which emit PM2.5 and black carbon at rates that are highly sensitive to combustion temperature and coal quality. The combined effect of near-zero ventilation and high emission rates produces daily-mean PM2.5 concentrations that routinely exceed 100 μg/m³ at peak winter, placing the city among the most polluted in the IQAir annual rankings for multiple consecutive years \cite{iqair_world_2024}. The WHO 24-hour PM2.5 guideline is 15 μg/m³ \cite{who_aqg_2021}; Almaty regularly exceeds this threshold by a factor of six or more during inversion events, and monthly means of 80.7 μg/m³ (December 2025) and 74.2 μg/m³ (January 2026) were recorded in the sensor network used in this study.

Until approximately 2023, the public data infrastructure for monitoring this problem was severely constrained. Almaty's official Kazhydromet network operated fewer than 10 reference-grade monitoring stations within a 682 km² urban area, a density that is inadequate for capturing the fine-scale spatial gradients characteristic of basin-trapped pollution. The emergence of community-scale low-cost sensor (LCS) networks has changed this situation substantially. The AirGradient open-hardware platform, combined with the OpenAQ data aggregation API \cite{openaq_platform}, has enabled a rapid proliferation of low-cost monitors across Almaty. As of April 2026, the network used in the present study comprised 192 sensors within a 25 km radius of the city centre (43.24°N, 76.89°E), of which 153 were AirGradient units concentrated in central and southern residential districts. This density, one sensor per roughly 3.5 km² on average, with higher density in the most densely populated zones, provides a city-scale spatial picture that was impossible to obtain from the reference network alone. Other Central Asian cities facing analogous inversion-driven pollution problems, including Shymkent (Kazakhstan), Karaganda (Kazakhstan), Bishkek (Kyrgyzstan), and Tashkent (Uzbekistan), share the same broad data-availability profile: sparse official monitoring, a growing LCS presence on OpenAQ, and essentially no published machine-learning forecasting work tailored to their specific geography.

Short-term PM2.5 forecasting has matured considerably in the European and North American contexts where reference-grade training targets and dense meteorological networks are available. Gradient-boosted tree models and long short-term memory networks trained on ground-based or satellite-derived features routinely achieve useful skill at 6--24 h horizons. Almaty sits outside this literature for two distinct reasons. The first is data access: the global platforms that publish PM2.5 forecasts, IQAir, AirVisual, BreezoMeter, operate proprietary pipelines whose model architectures, training data, and evaluation protocols are not in the public domain. The second is institutional capacity: Kazhydromet publishes hourly observations at its monitoring stations but does not provide a public forecast API or machine-readable hazard thresholds. The scientific literature contains no peer-reviewed study that trains a gradient-boosted model on Almaty's LCS network and evaluates it against a calibrated freely available baseline on a held-out winter test period. This gap is the direct motivation for the present work.

The significance of this work extends beyond the specific case of Almaty. Central Asia as a region carries one of the highest per-capita PM2.5 disease burdens globally, driven by a combination of arid dust sources, coal-heavy energy systems, and rapidly growing cities. A reproducible pipeline that depends only on an OpenAQ API key, a free Open-Meteo weather archive endpoint \cite{zippenfenig_openmeteo_2023}, and standard Python libraries deployable on free-tier cloud infrastructure represents a transferable template for any LCS-served city in the region. The parametric design of the pipeline - specifically, the fact that city centre coordinates, API query windows, and training configuration can be changed in a single configuration file, means that researchers or municipal agencies in Bishkek, Shymkent, or Tashkent could adapt the system to their local network within hours rather than months, without requiring proprietary data or institutional data-sharing agreements.

\subsection{Problem statement}\label{problem-statement}

The present research addresses the following gap: there is no open, reproducible short-term PM2.5 forecasting system built on the LCS network that exists in Almaty, and no published evaluation of how a locally trained model compares to the freely available CAMS global forecast for this specific geography. The scientific gap is the absence of empirical evidence quantifying how much predictive skill is gained by training a machine-learning model on Almaty's own sensor data relative to applying an off-the-shelf global chemistry-transport model.

Two subsidiary questions are embedded in this gap. First, given the limited size of the available LCS dataset (approximately 9,242 training hours across 562 days), it is unclear whether a non-linear gradient-boosted model or a regularised linear model provides superior accuracy, the answer has direct implications for model selection in analogous data-scarce urban environments elsewhere in Central Asia. Second, the spatial and temporal representativeness of a city-median derived from an unevenly distributed LCS network has not been quantified for Almaty: during sensor outage periods, the median may reflect fewer than 30 active stations rather than the nominal 192, and the consequences for forecast validity are unknown. Addressing these questions requires a complete evaluation pipeline operating on real, openly accessible data from Almaty's actual sensor network rather than on synthetic or reference-grade surrogate data. The present study provides that pipeline and reports the first quantitative answers to both questions on a 562-day window ending April 2026.

\subsection{Scope of the intelligent system}\label{scope-of-the-intelligent-system}

The present thesis realises the general programme described in its title - \emph{development of an intelligent system for investigating and solving ecological problems}, through a focused, high-impact instantiation: wintertime PM2.5 forecasting in Almaty, Kazakhstan. The choice of PM2.5 in an inversion-dominated mountain basin is not arbitrary. It represents the most acute ecological health hazard in the region, the most severe data gap relative to the harm it causes, and the most tractable starting point for an open, transferable intelligent system. The architecture developed here, automated multi-source ingestion, BLH-aware feature engineering, multi-horizon supervised models, and a zero-authentication public forecast interface, is explicitly designed to generalise across ecological monitoring tasks. The same pipeline can ingest NO₂ or O₃ rather than PM2.5; the same coordinate-configurable query layer can serve Bishkek or Tashkent rather than Almaty. The thesis thus delivers both a working system for a specific ecological problem and a reusable template for the broader class of problems named in its title.

\subsection{Aim and objectives}\label{aim-and-objectives}

\textbf{Aim.} Develop and evaluate an intelligent short-term PM2.5 forecasting system for Almaty that is built entirely on open data and code and is deployable on free-tier infrastructure.

\textbf{Objectives.} 1. Ingest and curate PM2.5 observations (OpenAQ LCS) and meteorological features (Open-Meteo, including boundary-layer height) for the 2024-10-01 → 2026-04-15 window into a reproducible SQLite store. 2. Engineer temporal and meteorological features tailored to Almaty's winter inversion regime. 3. Train and cross-validate an XGBoost forecaster for 6/12/24 h horizons and compare it against three baselines: persistence, linear regression, and CAMS. 4. Analyse error regimes, in particular, episodes of sudden onset and rapid clearing, and articulate the model's failure modes. 5. Ship a single-screen Streamlit web application that visualises current observations and the XGBoost forecast on a map of Almaty.

\subsection{Scientific novelty}\label{scientific-novelty}

The first contribution of the present study is the development of the first publicly released PM2.5 forecasting model calibrated specifically for Almaty's LCS network and inversion-dominated winter meteorology. Prior work on Almaty air quality is predominantly descriptive (source apportionment, trend analysis), or relies on satellite-derived estimates rather than ground-level sensor data. The present model is trained and evaluated on the 192-station OpenAQ network over a 562-day window ending April 2026, using a feature set that explicitly encodes BLH dynamics and inversion-onset signals not present in earlier local work.

The second contribution is the reproducible open architecture. All inputs are openly licensed (CC BY 4.0, no paywalled data); the full pipeline is a single Python repository with no cloud dependencies. Comparable published systems - including those evaluated by \cite{enebish2021ulaanbaatar, menares2021forecasting, Ayus2023Comparison}, typically rely on proprietary ground truth, restricted API access, or institutional data agreements that prevent reproduction or geographic transfer. The present design explicitly eliminates these barriers, making the pipeline portable to Shymkent, Karaganda, and other Central Asian cities with analogous LCS-on-OpenAQ data profiles.

The third contribution is a methodologically significant result regarding model-complexity trade-offs. Ridge regression outperforms XGBoost at the 12 and 24-hour forecast horizons, by 9.8\% in RMSE at +12 h (Ridge: 25.85 vs. XGBoost: 28.67 μg/m³) and by 13.5\% at +24 h (Ridge: 27.24 vs.~XGBoost: 31.48 μg/m³). This result is observed on a training set of 9,242 hours with 27 features, a scale at which, on the basis of published benchmarks, non-linear ensemble methods would normally be expected to retain their advantage. The finding contributes to the evidence base on model-complexity trade-offs for air quality forecasting problems with under 10,000 training samples. The most closely related published comparisons \cite{enebish2021ulaanbaatar, menares2021forecasting} report similar model-size reversals in settings where the target variable carries substantial measurement noise, which is consistent with the known hygroscopic overestimation bias in uncalibrated LCS readings.

This reversal also has an immediate practical consequence for system design. Because Ridge regression requires no tree construction, its inference time on a single row of features is negligible compared to XGBoost, and its serialised model file is several orders of magnitude smaller. In the context of a free-tier Streamlit deployment where memory and compute are constrained, a Ridge model at the 12 and 24-hour horizons is not merely equally accurate but computationally preferable. This alignment of statistical and engineering considerations strengthens the case for a multi-model operational system over a single-model one.

\subsection{Object, subject, and methods}\label{object-subject-and-methods}

\begin{itemize}
\tightlist
\item
  \textbf{Object of study:} short-term PM2.5 concentration dynamics in Almaty.
\item
  \textbf{Subject of study:} machine-learning models for hourly PM2.5 forecasting on LCS-based observations in inversion-dominated mountain-basin conditions.
\item
  \textbf{Methods:} data engineering on openly available APIs; supervised learning with gradient-boosted trees; time-series cross-validation; baseline comparison; error-regime analysis.
\end{itemize}

\subsection{Practical significance}\label{practical-significance}

The locally trained XGBoost model achieves RMSE = 23.44 μg/m³ and R² = 0.545 at the 6-hour forecast horizon on a held-out winter test set of 2,311 hours (December 2025 -- April 2026), outperforming the CAMS global chemistry-transport model by 41.1\% in RMSE at the same horizon. At the 12-hour horizon, Ridge regression produces RMSE = 25.85 μg/m³ and R² = 0.446, outperforming CAMS by approximately 28\% in RMSE. At 24 hours, Ridge achieves RMSE = 27.24 μg/m³ and R² = 0.383, while CAMS produces negative R² values at all three horizons (−0.314, −0.376, and −0.184 at 6, 12, and 24 h respectively), indicating that the global model provides no useful predictive signal for this geography. All locally trained models outperform CAMS by at least 16\% in RMSE across every horizon tested.

The practical deployment of these results takes the form of a Streamlit web application that provides a continuously updated PM2.5 forecast for Almaty at no cost to the end user. The application refreshes on a 30--60 second cycle, requires no authentication, and can be accessed from any device with a browser, making it suitable for use by schools, hospitals, municipal public health offices, and community organisations that lack technical staff or data infrastructure. The forecast uncertainty is communicated as a confidence band of approximately ±17 μg/m³ at the 6-hour horizon (±1 MAE), which provides a directly operationally useful range: an advisory issued when the lower bound of the forecast band crosses the WHO 24-hour guideline of 15 μg/m³ will capture genuine deterioration events at a false-positive rate acceptable for public communication purposes. The map view displays the current city-median PM2.5 alongside the three horizon forecasts, and the colour scale is calibrated to the Kazakhstani national AQI categories so that non-specialist users can interpret the displayed values without reference to μg/m³ numerics. The complete pipeline, data ingest, feature engineering, model training, and application, is implemented in a single open Python repository with no cloud dependencies, making it directly portable to other Central Asian cities.

\begin{quote}
Deployment URL: {[}to be added after Week 3 deployment{]}
\end{quote}

\subsection{Structure of the thesis}\label{structure-of-the-thesis}

Chapter 1 reviews the scientific and policy context in three parts: the meteorological and emission drivers of Almaty's PM2.5 episodes, including the role of temperature inversions and coal combustion; the field of low-cost sensor technology and the OpenAQ data ecosystem that makes the present dataset possible; and the methodological literature on gradient-boosted tree models and alternative approaches for short-term urban PM2.5 forecasting, with emphasis on studies conducted in inversion-dominated mountain or basin environments. Chapter 2 describes the data sources and system architecture in full reproducible detail: the SQLite database schema, the OpenAQ and Open-Meteo API ingestion procedures, the 27-feature matrix construction (including BLH delta, temporal lag, and calendar features), the train/test split, and the rationale for each non-obvious design choice. Chapter 3 reports the experimental results: the quantitative model comparison across all three forecast horizons and all four models, the feature importance analysis, the error-regime case studies on inversion onset and clearing events, and the spatial analysis of sensor-coverage effects during the October--November 2025 AirGradient firmware outage period. Chapter 4 discusses the implications of the results against the stated research questions, situates the findings within the published literature, quantifies the principal limitations, and proposes five specific directions for future work, including SHAP-based feature attribution, forecast BLH integration, and an ensemble of XGBoost and Ridge regression. An appendix provides the full model hyperparameter configurations, the API query specifications, and the data-quality filter thresholds applied during ingestion, so that the complete pipeline can be reproduced from the archived code and database without ambiguity.
